\documentclass[aspectratio=169,10pt]{beamer}
\makeatletter\def\input@path{{../estilo/}{estilo/}}\makeatother
\usepackage{devmobile}
\modulo{07}{Persistência Local}{Opções de armazenamento, DataStore, Room, relações, migrações e integração com a arquitetura}

\begin{document}
\capa

\begin{frame}{Objetivos do módulo}
  \begin{conceito}[Ao final deste módulo você será capaz de]
    \begin{itemize}
      \item Escolher a forma de armazenamento adequada a cada tipo de dado.
      \item Guardar preferências do usuário com \textbf{DataStore}.
      \item Modelar e consultar um banco SQLite com \textbf{Room}: entidades, DAOs e banco.
      \item Representar relações entre tabelas e converter tipos.
      \item Evoluir o esquema do banco com migrações sem perder dados.
      \item Integrar o banco ao repositório e expor dados reativos com \cod{Flow}.
    \end{itemize}
  \end{conceito}
\end{frame}

\sumario

% =============================================================================
\section{Opções de armazenamento}

\begin{frame}{Onde guardar cada coisa}
  \footnotesize
  \begin{tabularx}{\textwidth}{@{}l L L@{}}
    \toprule
    \cabecalho{Mecanismo} & \cabecalho{Indicado para} & \cabecalho{Exemplo no LabMonitor} \\
    \midrule
    \textbf{DataStore Preferences} & pares chave-valor pequenos & tema, unidade (°C/°F), limite de alerta \\
    \textbf{Proto DataStore} & objeto de configuração tipado & perfil do usuário em cache \\
    \textbf{Room (SQLite)} & dados estruturados, consultas, relações & equipamentos, bancadas, leituras \\
    \textbf{Arquivos internos} & arquivos privados do app & CSV exportado, logs \\
    \textbf{Cache} (\cod{cacheDir}) & dados descartáveis & miniaturas baixadas \\
    \textbf{MediaStore / seletor} & arquivos compartilhados com o usuário & fotos de equipamentos na galeria \\
    \bottomrule
  \end{tabularx}
  \vspace{4pt}
  \begin{columns}[T]
    \begin{column}{0.48\textwidth}
      \begin{conceito}[Armazenamento interno]
        Pasta privada do app (\cod{filesDir}). Outros apps não acessam, e ela é apagada na desinstalação.
      \end{conceito}
    \end{column}
    \begin{column}{0.48\textwidth}
      \begin{atencao}[SharedPreferences]
        A API antiga é síncrona e pode travar a UI. Em código novo, use DataStore.
      \end{atencao}
    \end{column}
  \end{columns}
\end{frame}

% =============================================================================
\section{DataStore}

\begin{frame}[fragile]{DataStore Preferences}
  \begin{columns}[T]
    \begin{column}{0.56\textwidth}
\begin{codigo}[title=AjustesRepository.kt]{Kotlin}
private val Context.ajustes by preferencesDataStore(
  name = "ajustes")

class AjustesRepository @Inject constructor(
  @ApplicationContext private val ctx: Context
) {
  private object Chaves {
    val LIMITE = doublePreferencesKey("limite_alerta")
    val FAHRENHEIT = booleanPreferencesKey("usar_f")
  }

  val limite: Flow<Double> = ctx.ajustes.data
    .catch { emit(emptyPreferences()) }  // erro de I/O
    .map { it[Chaves.LIMITE] ?: 30.0 }

  suspend fun definirLimite(valor: Double) {
    ctx.ajustes.edit { it[Chaves.LIMITE] = valor }
  }
}
\end{codigo}
    \end{column}
    \begin{column}{0.42\textwidth}
      \begin{conceito}[Características]
        \begin{itemize}
          \item Leitura como \cod{Flow}: a UI atualiza sozinha.
          \item Escrita transacional com \cod{edit}.
          \item Main-safe: I/O fora da thread principal.
        \end{itemize}
      \end{conceito}
      \begin{dica}[Dependência]
        \cod{androidx.datastore:}\\\cod{datastore-preferences}
      \end{dica}
      \begin{atencao}
        Uma instância por arquivo; listas grandes vão para o Room.
      \end{atencao}
    \end{column}
  \end{columns}
\end{frame}

% =============================================================================
\section{Room}

\begin{frame}{Arquitetura do Room}
  \begin{columns}[c]
    \begin{column}{0.52\textwidth}
      \centering
      \begin{tikzpicture}[node distance=6mm]
        \node[caixaescura,minimum width=4.6cm] (app) {Repositório / ViewModel};
        \node[caixaambarforte,minimum width=4.6cm,below=of app] (db) {\cod{@Database} AppDatabase};
        \node[caixa,minimum width=2.1cm,below left=6mm and -21mm of db] (dao) {\cod{@Dao}\\ BancadaDao};
        \node[caixa,minimum width=2.1cm,below right=6mm and -21mm of db] (ent) {\cod{@Entity}\\ Bancada};
        \node[caixa,minimum width=4.6cm,below=17mm of db,fill=cinzaclaro] (sql) {SQLite (arquivo \arq{labmonitor.db})};
        \draw[seta] (app) -- node[rotulo,right]{obtém DAOs} (db);
        \draw[seta] (db) -- (dao);
        \draw[seta] (db) -- (ent);
        \draw[seta] (dao) -- node[rotulo,left]{SQL gerado} (sql);
        \draw[setatrac] (ent) -- node[rotulo,right]{tabela} (sql);
      \end{tikzpicture}
    \end{column}
    \begin{column}{0.46\textwidth}
      \small
      \termo{Room} é uma camada sobre o SQLite que:
      \begin{itemize}
        \item verifica o SQL \textbf{em tempo de compilação};
        \item converte linhas em objetos Kotlin automaticamente;
        \item devolve \cod{Flow} que reemite quando a tabela muda;
        \item suporta coroutines e migrações;
        \item funciona também em Kotlin Multiplatform.
      \end{itemize}
      \begin{dica}[Dependências]
        \cod{room-runtime}, \cod{room-ktx} e \cod{ksp(room-compiler)}.
      \end{dica}
    \end{column}
  \end{columns}
\end{frame}

\begin{frame}[fragile]{Entidades}
  \begin{columns}[T]
    \begin{column}{0.54\textwidth}
\begin{codigo}[title=Entidades.kt]{Kotlin}
@Entity(tableName = "bancadas")
data class BancadaEntity(
  @PrimaryKey val id: Int,
  val nome: String,
  val sala: String,
  @ColumnInfo(name = "limite_c") val limite: Double = 30.0
)

@Entity(
  tableName = "leituras",
  foreignKeys = [ForeignKey(
    entity = BancadaEntity::class,
    parentColumns = ["id"], childColumns = ["bancadaId"],
    onDelete = ForeignKey.CASCADE)],
  indices = [Index("bancadaId")]
)
data class LeituraEntity(
  @PrimaryKey(autoGenerate = true) val id: Long = 0,
  val bancadaId: Int,
  val valor: Double,
  val instante: Instant
)
\end{codigo}
    \end{column}
    \begin{column}{0.44\textwidth}
      \centering
      \begin{tikzpicture}
        \node[caixa,font=\tiny,align=left,text width=1.9cm] (b) {\textbf{bancadas}\\ \underline{id} INTEGER\\ nome TEXT\\ sala TEXT\\ limite\_c REAL};
        \node[caixaambar,font=\tiny,align=left,text width=2.3cm,right=7mm of b] (l) {\textbf{leituras}\\ \underline{id} INTEGER\\ bancadaId INTEGER (FK)\\ valor REAL\\ instante INTEGER};
        \draw[seta] (l.west) -- node[rotulo,above]{N:1} (b.east);
      \end{tikzpicture}
      \begin{conceito}[Anotações]
        \cod{@Entity} vira tabela, \cod{@PrimaryKey} define a chave e \cod{@ColumnInfo} renomeia colunas. \cod{ForeignKey} garante a integridade e o \cod{CASCADE} apaga as leituras junto com a bancada.
      \end{conceito}
      \begin{dica}
        Crie \cod{Index} nas colunas usadas em \cod{WHERE} e em chaves estrangeiras.
      \end{dica}
    \end{column}
  \end{columns}
\end{frame}

\begin{frame}[fragile]{DAO: \emph{Data Access Object}}
  \begin{columns}[T]
    \begin{column}{0.56\textwidth}
\begin{codigo}[title=BancadaDao.kt]{Kotlin}
@Dao
interface BancadaDao {
  @Query("SELECT * FROM bancadas ORDER BY nome")
  fun observarTodas(): Flow<List<BancadaEntity>>

  @Query("SELECT * FROM bancadas WHERE id = :id")
  suspend fun buscar(id: Int): BancadaEntity?

  @Upsert                         // insere ou atualiza
  suspend fun salvar(vararg b: BancadaEntity)

  @Delete suspend fun remover(b: BancadaEntity)
  @Query("""SELECT AVG(valor) FROM leituras
    WHERE bancadaId = :id AND instante >= :desde""")
  fun media(id: Int, desde: Instant): Flow<Double?>
  @Transaction
  suspend fun substituirTodas(lista: List<BancadaEntity>) {
    apagarTodas(); salvar(*lista.toTypedArray())
  }
  @Query("DELETE FROM bancadas") suspend fun apagarTodas()
}
\end{codigo}
    \end{column}
    \begin{column}{0.42\textwidth}
      \scriptsize
      \begin{tabularx}{\linewidth}{@{}l L@{}}
        \toprule
        \cabecalho{Retorno} & \cabecalho{Comportamento} \\
        \midrule
        \cod{Flow<T>} & observa: reemite a cada mudança \\
        \cod{suspend fun} & executa uma vez, fora da UI \\
        \bottomrule
      \end{tabularx}
      \begin{conceito}[Parâmetros]
        \cod{:id} na consulta liga-se ao parâmetro \cod{id} da função. O Room valida nomes de tabelas e colunas na compilação.
      \end{conceito}
      \begin{dica}[\cod{@Transaction}]
        Agrupa operações: ou todas acontecem, ou nenhuma. Evita que a UI veja a tabela vazia no meio da troca.
      \end{dica}
    \end{column}
  \end{columns}
\end{frame}

\begin{frame}[fragile]{Banco de dados e conversores de tipo}
  \begin{columns}[T]
    \begin{column}{0.54\textwidth}
\begin{codigo}[title=AppDatabase.kt]{Kotlin}
@Database(
  entities = [BancadaEntity::class,
              LeituraEntity::class],
  version = 1,
  exportSchema = true
)
@TypeConverters(Conversores::class)
abstract class AppDatabase : RoomDatabase() {
  abstract fun bancadaDao(): BancadaDao
  abstract fun leituraDao(): LeituraDao
}

class Conversores {
  @TypeConverter
  fun deInstant(i: Instant?): Long? = i?.toEpochMilli()
  @TypeConverter
  fun paraInstant(ms: Long?): Instant? =
    ms?.let(Instant::ofEpochMilli)
}
\end{codigo}
    \end{column}
    \begin{column}{0.44\textwidth}
      \begin{conceito}[Conversores]
        O SQLite só conhece inteiro, real, texto e blob. \cod{@TypeConverter} ensina o Room a guardar \cod{Instant}, \cod{enum} e outros tipos.
      \end{conceito}
      \begin{dica}[Criação via Hilt (Módulo 06)]
        \cod{Room.databaseBuilder(...)} em um \cod{@Provides @Singleton}: uma instância única para o app inteiro.
      \end{dica}
      \begin{dica}
        \cod{createFromAsset("inicial.db")} cria o banco já preenchido (ex.: inventário do laboratório).
      \end{dica}
    \end{column}
  \end{columns}
\end{frame}

\begin{frame}[fragile]{Relações entre tabelas}
  \begin{columns}[T]
    \begin{column}{0.54\textwidth}
\begin{codigo}[title=Um-para-muitos]{Kotlin}
data class BancadaComEquipamentos(
  @Embedded val bancada: BancadaEntity,
  @Relation(
    parentColumn = "id",
    entityColumn = "bancadaId"
  )
  val equipamentos: List<EquipamentoEntity>
)

@Dao
interface BancadaDao {
  @Transaction
  @Query("SELECT * FROM bancadas WHERE id = :id")
  fun comEquipamentos(id: Int):
    Flow<BancadaComEquipamentos?>
}
\end{codigo}
    \end{column}
    \begin{column}{0.44\textwidth}
      \begin{conceito}[\cod{@Embedded} + \cod{@Relation}]
        O Room faz duas consultas e monta o objeto composto: a bancada e a lista dos seus equipamentos.
      \end{conceito}
      \begin{dica}[Muitos-para-muitos]
        Crie uma tabela de junção (ex.: \cod{ReservaUsuario}) e use \cod{@Relation(associateBy = Junction(...))}.
      \end{dica}
      \begin{atencao}
        Use \cod{@Transaction} em consultas com \cod{@Relation} para que as duas consultas leiam um estado consistente.
      \end{atencao}
    \end{column}
  \end{columns}
\end{frame}

\begin{frame}[fragile]{Migrações: evoluindo o esquema}
  \begin{columns}[T]
    \begin{column}{0.54\textwidth}
\begin{codigo}[title=De v1 para v3]{Kotlin}
@Database(entities = [...], version = 3,
  autoMigrations = [AutoMigration(from = 1, to = 2)])
abstract class AppDatabase : RoomDatabase()

// mudanças complexas: migração manual
val MIGRACAO_2_3 = object : Migration(2, 3) {
  override fun migrate(db: SupportSQLiteDatabase) {
    db.execSQL("ALTER TABLE bancadas " +
      "ADD COLUMN responsavel TEXT")
    db.execSQL("CREATE INDEX idx_leitura_inst " +
      "ON leituras(instante)")
  }
}

Room.databaseBuilder(ctx, AppDatabase::class.java, "lab.db")
  .addMigrations(MIGRACAO_2_3)
  .build()
\end{codigo}
    \end{column}
    \begin{column}{0.44\textwidth}
      \begin{atencao}[Regra]
        Mudou uma entidade? \textbf{Incremente a versão} e forneça a migração, senão o app trava ao abrir o banco antigo.
      \end{atencao}
      \begin{conceito}[\cod{exportSchema}]
        Gera um JSON do esquema de cada versão. Versione esses arquivos: são a base das \emph{auto migrations} e dos testes de migração.
      \end{conceito}
      \begin{erro}
        \cod{fallbackToDestructiveMigration()} apaga os dados do usuário. Aceitável só em protótipos ou caches.
      \end{erro}
    \end{column}
  \end{columns}
\end{frame}

% =============================================================================
\section{Integração com a arquitetura}

\begin{frame}{Fluxo reativo de ponta a ponta}
  \begin{center}
  \begin{tikzpicture}[node distance=5mm]
    \node[caixa,fill=cinzaclaro] (sql) {SQLite};
    \node[caixa,right=of sql] (dao) {DAO\\ \tiny \cod{Flow<List<Entity>>}};
    \node[caixapetroleo,right=of dao] (repo) {Repositório\\ \tiny \cod{Flow<List<Modelo>>}};
    \node[caixaescura,right=of repo] (vm) {ViewModel\\ \tiny \cod{StateFlow<UiState>}};
    \node[caixaambarforte,right=of vm] (ui) {UI\\ \tiny recompõe};
    \draw[seta] (sql) -- (dao); \draw[seta] (dao) -- (repo); \draw[seta] (repo) -- (vm); \draw[seta] (vm) -- (ui);
    \draw[setaambar,rounded corners] (ui.south) -- ++(0,-7mm) -| node[rotulo,pos=0.2,below]{evento: \cod{onSalvar()} $\rightarrow$ \cod{viewModelScope.launch \{ repo.salvar() \}} $\rightarrow$ \cod{dao.salvar()}} (sql.south);
  \end{tikzpicture}
  \end{center}
  \vspace{2pt}
  \begin{columns}[T]
    \begin{column}{0.48\textwidth}
      \begin{conceito}[A mágica]
        Depois que o DAO grava, o Room \textbf{notifica} todos os \cod{Flow} que leem aquela tabela. A tela se atualiza sem nenhum código extra de ``recarregar''.
      \end{conceito}
    \end{column}
    \begin{column}{0.48\textwidth}
      \begin{dica}[Inspecionando]
        \textbf{App Inspection $\rightarrow$ Database Inspector} mostra as tabelas do app em execução, permite rodar SQL e atualiza ao vivo.
      \end{dica}
    \end{column}
  \end{columns}
\end{frame}

\begin{frame}[fragile]{Guardando leituras de sensores}
  \begin{columns}[T]
    \begin{column}{0.56\textwidth}
\begin{codigo}[title=LeituraRepository.kt]{Kotlin}
class LeituraRepository @Inject constructor(
  private val dao: LeituraDao
) {
  fun ultimas(bancada: Int, n: Int = 120) =
    dao.ultimas(bancada, n)          // Flow

  suspend fun registrar(l: Leitura) =
    dao.inserir(l.paraEntidade())

  // mantém o banco pequeno: apaga o que é antigo
  suspend fun limpar(dias: Long = 30) =
    dao.apagarAntesDe(
      Instant.now().minus(dias, ChronoUnit.DAYS))
}

@Query("""SELECT * FROM leituras WHERE bancadaId = :b
  ORDER BY instante DESC LIMIT :n""")
fun ultimas(b: Int, n: Int): Flow<List<LeituraEntity>>
\end{codigo}
    \end{column}
    \begin{column}{0.42\textwidth}
      \begin{atencao}[Volume de dados]
        Um sensor a 1\,Hz gera 86\,400 linhas por dia. Planeje:
        \begin{itemize}
          \item retenção (apagar o que é antigo);
          \item agregação (médias por minuto);
          \item consultas com \cod{LIMIT} e índices.
        \end{itemize}
      \end{atencao}
      \begin{dica}
        Para inserir muitas leituras, agrupe em lote (\cod{@Insert} com \cod{List}) em vez de uma transação por leitura.
      \end{dica}
    \end{column}
  \end{columns}
\end{frame}

\begin{frame}{Boas práticas e segurança}
  \begin{columns}[T]
    \begin{column}{0.48\textwidth}
      \begin{conceito}[Boas práticas]
        \begin{itemize}
          \item Nunca acesse o banco na thread principal (o Room bloqueia por padrão).
          \item Uma instância de \cod{AppDatabase} por processo.
          \item Teste as migrações com \cod{MigrationTestHelper}.
          \item Banco em memória (\cod{inMemoryDatabaseBuilder}) nos testes (Módulo 11).
        \end{itemize}
      \end{conceito}
    \end{column}
    \begin{column}{0.48\textwidth}
      \begin{atencao}[Dados sensíveis]
        \begin{itemize}
          \item Não guarde senhas. Guarde \emph{tokens} e prefira o \textbf{Android Keystore} para chaves criptográficas.
          \item O backup automático do Android copia arquivos do app para a nuvem do usuário; exclua o que não deve sair do aparelho com \cod{dataExtractionRules}.
          \item Trate dados pessoais conforme a \textbf{LGPD} (Módulo 12).
        \end{itemize}
      \end{atencao}
    \end{column}
  \end{columns}
\end{frame}

% =============================================================================
\section{Fechamento}

\begin{frame}{Resumo}
  \begin{resumo}
    \begin{itemize}
      \item Preferências pequenas no DataStore; dados estruturados no Room; arquivos em \cod{filesDir}.
      \item Room = \cod{@Entity} (tabelas) + \cod{@Dao} (consultas) + \cod{@Database} (o banco).
      \item DAOs devolvem \cod{Flow} para observar e usam \cod{suspend} para operações pontuais.
      \item \cod{@TypeConverter} para tipos especiais; \cod{@Relation} para relações.
      \item Toda mudança de esquema exige nova versão e migração.
      \item Com repositório + \cod{Flow}, gravar no banco atualiza a tela automaticamente.
    \end{itemize}
  \end{resumo}
\end{frame}

\begin{frame}{Exercícios}
  \begin{exercicio}[Exercício 1 --- DataStore]
    Crie a tela de Ajustes com limite de alerta (\cod{Slider}), unidade de temperatura (\cod{Switch}) e intervalo de amostragem. Todos os valores devem persistir.
  \end{exercicio}
  \begin{exercicio}[Exercício 2 --- Consultas]
    Escreva consultas DAO para: as 5 bancadas com maior média de temperatura nas últimas 24\,h; a contagem de leituras acima do limite por bancada; leituras entre duas datas.
  \end{exercicio}
  \begin{exercicio}[Exercício 3 --- Migração]
    Adicione a coluna \cod{responsavel} à tabela de bancadas e escreva a migração 1$\rightarrow$2. Confirme no Database Inspector que os dados antigos foram preservados.
  \end{exercicio}
  \begin{projeto}[LabMonitor --- Etapa 7]
    Substitua os repositórios \emph{fake} por Room: bancadas, equipamentos e leituras do simulador. Guarde o limite de alerta no DataStore e implemente a limpeza de leituras com mais de 30 dias.
  \end{projeto}
\end{frame}

\begin{frame}{Para aprofundar}
  \small
  \begin{itemize}
    \item Visão geral de armazenamento: \url{https://developer.android.com/training/data-storage}
    \item DataStore: \url{https://developer.android.com/topic/libraries/architecture/datastore}
    \item Room: \url{https://developer.android.com/training/data-storage/room}
    \item Relações no Room: \url{https://developer.android.com/training/data-storage/room/relationships}
    \item Migrações: \url{https://developer.android.com/training/data-storage/room/migrating-db-versions}
    \item Backup automático: \url{https://developer.android.com/identity/data/autobackup}
  \end{itemize}
\end{frame}

\encerramento{Módulo 08 --- Rede e consumo de APIs: HTTP, REST, Retrofit, serialização, erros e estratégia offline-first.}

\end{document}
